\section{Methods}
\label{methods}

\subsection{Mathematical Formulation}
To contextualise the algorithm development carried out, the mathematical background has been included. The following section is divided into two main areas:

\begin{enumerate}
    \item \nameref{sec:era-proof}
    \item \nameref{sec:eigendecomposition}
\end{enumerate}

The core ERA derivation (Section~\ref{sec:era-proof}) follows \cite{dafx2026}; the eigendecomposition and modal-form realisation. The author makes no claim of originality or contribution of the mathematics in this particular section, and instead has been included in full to contextualise the additional work carried out in algorithm development and parallel filter construction. There are only small adaptions applied to the eigendecomposition description, which are largely inconsequential to the results.

\subsubsection{Eigensystem Realisation}
\label{sec:era-proof}

The full treatment and derivation of the ERA can be found in \cite{juang-1994-book}.

A system is exposed to a unit impulse input signal $u[k] = \delta[k]$ with initial condition of $x[0] = 0$. The impulse response of the system, $y[k]$ (also defined as Markov parameters) are as follows:
\begin{subequations}
\label{eq:markov}
\begin{align}
    \label{eq:markov1}
    y[0] &= d,\\
    \label{eq:markov2}
    y[1] &= \mathbf{c}^\top \mathbf{b},\\
    \label{eq:markov3}
    y[k] &= \mathbf{c}^\top \mathbf{A}^{k-1} \mathbf{b}, \quad k \geq 2
\end{align}
\end{subequations}
A Hankel matrix is then defined based on the Markov parameters:
\begin{align}
\label{eq:hankel}
    \mathbf{H}_0 = 
    \begin{bmatrix} 
        y[1] & y[2] & \dots & y[s] \\
        y[2] & y[3] & \dots & y[s+1] \\
        \vdots & \vdots & \ddots & \vdots \\
        y[r] & y[r+1] & \dots & y[r+s-1]
    \end{bmatrix},
\end{align}
alongside the shifted Hankel matrix:
\begin{align}
\label{eq:shift-hankel}
    \mathbf{H}_1 = 
    \begin{bmatrix} 
        y[2] & y[3] & \dots & y[s+1] \\
        y[3] & y[4] & \dots & y[s+2] \\
        \vdots & \vdots & \ddots & \vdots \\
        y[r+1] & y[r+2] & \dots & y[r+s]
    \end{bmatrix},
\end{align}
where $r,s \in \mathbb{N}$. The values of $s$ and $r$ are set depending on the application. If they are too small, the system dynamics will not be appropriately captured, but if too large, could introduce noise sensitivity and numerical ill-conditioning.

When the system is free from noise, the Hankel matrix can be factorised by making substitutions from \eqref{eq:markov}:
\begin{align}
\label{eq:factorised}
    \mathbf{H}_0 = 
    \begin{bmatrix}
    \mathbf{c}^{\top}, & \mathbf{c}^{\top} \mathbf{A}, & \dots, & \mathbf{c}^{\top} \mathbf{A}^{r-1}
    \end{bmatrix}^{\top}
    \begin{bmatrix}
    \mathbf{b}, & \mathbf{A} \mathbf{b}, & \dots, & \mathbf{A}^{s-1} \mathbf{b}
    \end{bmatrix}
    = \mathcal{O}_r \mathcal{C}_s,
\end{align}
where $\mathcal{O}_r$ and $\mathcal{C}_s$ are the extended observability and controllability matrices. Given that $r$ and $s$ are large enough, the rank of $\mathbf{H}_0$ is a minimal realisation of the system.

In other words, given that enough information has been captured from the impulse response, a sufficiently large number of rows and columns utilised in the Hankel matrix will encode the entirety of the system dynamics. In a real world scenario, measurement noise will be captured in the Hankel matrix, thus the dominant modes of the system are mixed with modes corresponding to the noise. Separation of the noise modes of the system is carried out via the Singular Value Decomposition (SVD):
\begin{align}
\label{eq:svd}
    \mathbf{H}_0 = \mathbf{U \Sigma V}^\top
\end{align}
By truncating the number $m$ of singular values so that only the dominant modes are retained, a rank-$m$ approximation yields:
\begin{align}
\label{eq:svdm}
    \mathbf{H}_0 \approx \mathbf{U}_m \mathbf{\Sigma}_m \mathbf{V}_{m}^{\top}
\end{align}
where $\Sigma_m =$ diag$(\sigma_1,\sigma_2, \dots, \sigma_m)$ are the largest singular value entries. $U_m$ and $V_m$ are the truncated left and right singular vectors respectively.

The state space matrices are composed as follows:
\begin{subequations}
\begin{align}
\label{eq:ss}
\hat{\mathbf{A}} &= \mathbf{\Sigma}_{m}^{-1/2} \mathbf{U}^{\top}_m \mathbf{H}_1 \mathbf{V}_m \mathbf{\Sigma}_{m}^{-1/2}, \\
\hat{\mathbf{b}} &= \mathbf{\Sigma}_{m}^{1/2} \mathbf{V}_m^{\top} \mathbf{e}_{1}^{(s)}, \\
\hat{\mathbf{c}} &= \mathbf{\Sigma}_{m}^{1/2} \mathbf{U}_{m}^{\top} \mathbf{e}_{1}^{(r)}, \\
\hat{d} &= y[0],
\end{align}
\end{subequations}

where $\mathbf{e}_{1}^{(\ell)} = [1, 0, ..., 0]^{\top}$ is a canonical basis vector with dimension $\ell$. The reconstruction of the approximate states and output can now be represented in the canonical form:
\begin{align}
\label{eq:output}
\begin{cases}
    \hat{\mathbf{x}}[k+1] = \hat{\mathbf{A}} \hat{\mathbf{x}} [k] + \hat{\mathbf{b}} u[k], \\
    \hat{y} [k] = \hat{\mathbf{c}}^{\top} \hat{\mathbf{x}} [k] + \hat{d} u[k].
\end{cases}
\end{align}
The shifted Hankel matrix $\mathbf{H}_1 = \mathcal{O}_r \mathbf{A} \mathbf{C}_s$ allows the recovery of the state transition matrix $\hat{\mathbf{A}}$. When the rank of $\hat{\mathbf{A}}$ matches the true system order, the result is the minimum realisation of the system. This provides a very explicit trade-off between the computational complexity and accuracy of the impulse response reconstruction.

\subsubsection{Eigendecomposition}
\label{sec:eigendecomposition}
Diagonalisation of the state space matrix $\hat{\mathbf{A}}$ and associated transforms on $\hat{\mathbf{b}}$ and $\hat{\mathbf{c}}$ allow the system to be represented in terms of linearly independent modes via eigendecomposition $\hat{\mathbf{A}} = \mathbf{Q} \Lambda \mathbf{Q}^{-1}$:
\begin{equation}
\label{eq:coord}
    \tilde{\mathbf{x}}       \leftarrow \mathbf{Q}^{-1}\mathbf{\hat{x}},
    \quad \tilde{\mathbf{A}} \leftarrow \Lambda,
    \quad \tilde{\mathbf{b}} \leftarrow \mathbf{Q}^{-1} \hat{\mathbf{b}},
    \quad \tilde{\mathbf{c}} \leftarrow \mathbf{Q}^{\top} \hat{\mathbf{c}}
\end{equation}
Where $\mathbf{Q}$ is the matrix of eigenvectors of the system matrix $\mathbf{A}$ and $\mathbf{\Lambda}$ is the matrix of eigenvalues. As the eigenvalues of the system are invariant under transformation, they allow the expression of each mode of the system in a linearly independent basis.
The eigenvalues of the system contain the frequency and damping factor information, with amplitude and phase residing in the residues. As the eigenvalues are in discrete time, they are first transformed into the continuous time domain via the following relations [INSERT ERA PPT REF HERE]:
\begin{equation}
\label{eq:era-parameters}
    \lambda_{c,i} = \frac{\ln \lambda_{d,i}}{\tau}, \quad 
    f_{i} = \frac{\abs{\Im(\lambda_{c,i})}}{2\pi}, \quad 
    \zeta_{i} = \frac{-\Re(\lambda_{c,i})}{\abs{\lambda_{c,i}}}, \quad
    A_{i} = |c_{i}b_{i}|, \quad
    \phi_{i} = \arg{(c_{i}b_{i})}.
\end{equation}
where $\lambda_{c,i}$ is the continuous time eigenvalue at index $i$, $\lambda_{d,i}$ is the associated discrete time eigenvalue, $\tau$ is the timestep in seconds, $f_{i}$ is the associated continuous time frequency and $\zeta_{i}$ is the continuous time damping factor.
After making the co-ordinate substitutions for \eqref{eq:markov}, the output is represented as follows:
\begin{align}
\label{eq:markov-eig}
    \mathbf{\tilde{x}}[k + 1] = \mathbf{\Lambda} \mathbf{ \tilde{x}}[k] + \mathbf{Q^{-1} \hat{b}}u[k]\\
    \tilde{y}[k] = \left( \mathbf{Q}^{\top} \hat{\mathbf{c}} \right)^{\top} \tilde{\mathbf{x}}[k] + \hat{d}u[k]
\end{align}
By exploiting the diagonal nature of the eigenvalue matrix, a further co-ordinate transformation can be carried out by taking the diagonal of $\mathbf{\Lambda}$, $\tilde{\mathbf{A}}$ = diag$(\mathbf{\Lambda}) = [\lambda_1, \lambda_2, \dots, \lambda_m]^{\top}$. Utilising this property provides an efficient and succinct method for computation of the reconstructed impulse response and physical modes of the system as an expression of powers of $\tilde{\mathbf{A}}$:
\begin{align}
\label{eq:markov-again}
\begin{cases}
    \mathbf{\tilde{x}}[k + 1] =  \tilde{\mathbf{A}}^{k-1} \mathbf{Q}^{-1} \hat{\mathbf{b}}, \\
    \tilde{y}[k] =  \hat{\mathbf{c}}^{\top} \mathbf{Q} \tilde{\mathbf{A}}^{k-1}\mathbf{Q}^{-1}\hat{\mathbf{b}}, \quad k \geq 1
\end{cases}
\end{align}

\subsubsection{Modal Form Realisation}
\label{sec:modal-form}
A parallel IIR filter formulation requires the partial fraction expansion of the complex conjugates for each mode in the system to retrieve the corresponding filter coefficients. Consider the standard discrete time transfer function for a state space system [INSERT REF HERE]:
\begin{equation}
\label{eq:ss-tf}
    \mathbf{H}(z) = \mathbf{C}(z \mathbf{I} - \mathbf{A})^{-1} \mathbf{B} + \mathbf{D}.
\end{equation}
Making the substitution for the diagonalised form of the system from \eqref{eq:coord} yields the following explicit form for a single mode system:
\begin{equation}
\label{eq:modal}
    H(z) = \begin{bmatrix} c_1 & c_2 \end{bmatrix} \left( z \mathbf{I} - \begin{bmatrix} \lambda_1 & 0 \\ 0 & \lambda_2 \end{bmatrix} \right)^{-1}
    \begin{bmatrix} b_1 \\ b_2 \end{bmatrix} + d.
\end{equation}
It should be noted that $d$ in the context of ERA is simply a scalar term observed as the first sample of the signal. Expanding the transfer function explicitly will give rise to the parallel single-pole representation of the system:
\begin{equation}
\label{eq:one-pole}
    H(z) = \frac{c_1 b_1}{z - \lambda_1} + \frac{c_2 b_2}{z - \lambda_2} + d
\end{equation}
where $c_2 = \overline{c}_1$, $b_2 = \overline{b}_1$ and $\lambda_2 = \overline{\lambda}_1$.
Performing partial fraction expansion on \eqref{eq:one-pole} and multiplying by $z^{-2}$ gives rise to the parallel second order form which is necessary for the IIR filter implementation:
\begin{equation}
\label{eq:second-order}
    H(z) = \frac{(R_1 + R_2)z^{-1} - (R_1 \lambda_2 + R_2 \lambda_1)z^{-2}}{1 - (\lambda_1 + \lambda_2)z^{-1} + |\lambda_1|^{2}z^{-2}} + d,
\end{equation}

where $R_1$ is the residue $c_1 b_1$, and $R_2$ is the residue $c_2 b_2$, which as previously highlighted, is the complex conjugate of $R_1$.
By simplifying further by exploiting the conjugate symmetry and folding the free term $+d$ into the expression:
\begin{equation}
\label{eq:full-parallel}
    H(z) = \frac{d + \left[2\Rey(R_1) - 2d\Rey(\lambda_1)\right]z^{-1} + \left[d|\lambda_1|^2 - 2\Rey(R_1\lambda_2)\right]z^{-2}}{1 - 2\Rey(\lambda_1)z^{-1} + |\lambda_1|^2 z^{-2}}.
\end{equation}

Comparing this with the standard form of second order IIR filter, the coefficients can be represented as follows:
\begin{equation}
\label{eq:iir-filter}
H(z) = \frac{b_0 + b_1 z^{-1} + b_2 z^{-2}}{1 + a_1 z^{-1} + a_2 z^{-2}}
\end{equation}
where the coefficients for a single mode are:
\begin{subequations}
\label{eq:biquad-coeffs}
\begin{align}
    b_0 &= d, \\
    b_1 &= 2\Rey(R_1) - 2d\Rey(\lambda_1), \\
    b_2 &= d|\lambda_1|^2 - 2\Rey(R_1\lambda_2), \\
    a_1 &= -2\Rey(\lambda_1), \\
    a_2 &= |\lambda_1|^2.
\end{align}
\end{subequations}
It is obvious from this description that systems with a higher order of modes, $d$ can be folded into one filter section, thus the numerator coefficients reduce to,
\begin{subequations}
\label{eq:biquad-coeffs-2}
\begin{align}
    b_0 &= 0, \\
    b_1 &= 2\Rey(R_1), \\
    b_2 &= -2\Rey(R_1\lambda_2), \\
\end{align}
\end{subequations}
where $a_1$ and $a_2$ remain unchanged.

\subsubsection{FDM Modal Form}
In \cite{oneill}, the Impulse Invariant Method (IIM) \cite{dsp-book} is utilised to convert the continuous time parameters produced from the FDM method to the IIR filter coefficients for each mode:
\begin{equation}
\label{iim-tf}
H(z) = A_{i} \frac{\sin(\phi_{i}) + R_{i} \left[ \cos(\phi_{i})S_{i} - \sin(\phi)C_{i}\right]z^{-1}}{1 - 2R_{i}C_{i}z^{-1} + R_{i}^{2}z^{-2}}
\end{equation}
where
\begin{equation}
\label{iim-coeff}
R_{i} = e^{-\alpha_{i}T}, \quad S_{i} = \sin(\omega_{i}T), \quad C_{i} = \cos(\omega_{i}T).
\end{equation}
These are then mapped to the second order filter coefficients as follows:
\begin{subequations}
\label{iim-iir-coeffs}
\begin{align}
    b_{0,i} &= A_{i}\sin(\phi_{i}), \\
    b_{1,i} &= A_{i}R_{i} \left[ \cos(\phi_{i})S_{i} - \sin(\phi)C_{i}\right], \\
    a_{1,i} &= -2R_{i}C_{i}, \\
    a_{2,i} &= R_{i}^2.
\end{align}
\end{subequations}
This formulation offers an advantage of the ERA method over the FDM, in that ERA derived poles and residues of the system are already in discrete time, the IIM can be bypassed completely; the poles and residues can be converted directly to filter coefficients.

\subsection{Framework Overview}
This section details major design decisions in the building of the framework. Detail is added to specific stages that arise from the mismatch between mathematical theory and computational implementation, or where the mathematical background fails to provide a satisfactory explanation.

\subsubsection{Hankel Dimensioning}
Hankel matrix dimensioning is carried out automatically when no inputs are provided. The Hankel dimensions are set as $\text{floor}(nSamples / 2) - 1$, creating a square matrix with the maximum number of samples available. Overdetermination of the Hankel Matrix size can cause too much of the noise of the system to be captured which has a knock on effect on subsequent steps of the algorithm. Additionally, underdetermination of the Hankel Matrix may cause the system to miss some physical modes due to insufficient data to accurately capture it's contribution, both points are raised in \cite{dafx2026}. To mitigate this, signals are recommended to be pre-truncated and filtered to reduce the degree of noise in the signal, allowing one dimension rule to be applied to each signal under investigation.

\subsubsection{SVD Truncation}
The SVD truncation method will vary dependant upon the value provided to the function. For values greater than 1, the SVD truncation routine will retain $2m$ singular values, where $m$ is the requested number of modes retained. If the requested number of modes is greater than the rank of the $\Sigma$ matrix, the nearest even number of modes is retained.
If the provided value is $0<m<1$, the tolerance will act as a cutoff level on the normalised singular values. The maximum even number of modes above the cutoff will be retained. In an earlier stage of the project, it was assumed that the SVD truncation would be useful for noise rejection within the signal. Whilst this is a byproduct of separating the lower order modes, no simple method was identified for separation of the signal subspace from noise subspace with this method, thus two methods of truncation were reproduced, similar in nature to \cite{pyyeti}.
\begin{figure}
\centering
\includegraphics{figures/out/svd-comparison.pdf}
\caption{\textnormal Comparison of SVD truncation methods, number of requested modes $m=5$ (left) and tolerance threshold set to 0.1 (right)}
\label{svd}
\end{figure}

\subsubsection{Modal Parameter Extraction}
Following SVD truncation, the state space models are constructed, according to \eqref{eq:ss}. As there are no deviations from the the mathematics, there are no design considerations to discuss. Modal parameter extraction contains the majority of the interesting discussion as to how pole recovery is treated.
Diagonalisation is carried out according to \eqref{eq:markov-again}, and the subsequent eigenvalues are sorted according to their imaginary components, similar to \cite{pyyeti}.
Once the modal parameters frequency, phase, amplitude and damping factor are extracted according to \eqref{eq:era-parameters}. A set of togglable filters are then applied to the system to remove unphysical modes. Filters are built as binary masks, and applied in several stages.
\begin{itemize}
    \item Stage 1 Mask:
    \begin{itemize}
        \item Unstable pole filter: poles with magnitudes $>=1$ are removed. These poles are unphysical and correspond to a growing exponential, which is physcially impossible.
        \item Zero frequency filter: poles with a zero frequency component (pure exponential decay) are removed. These poles violate the assumption of the mathematical framework explored in Section~\ref{sec:modal-form}, offsetting the conjugate pair indexing necessary to build the filter coefficients.
    \end{itemize}
    \item Stage 2 Mask:
    \begin{itemize}
        \item Conjugate pole verification: The pole array is reshaped into a $2\times\frac{p}{2}$ matrix, where $p$ is the number of retained poles. The array is then scanned to ensure that if one pole has been flagged for removal by the stage 1 mask, the associated conjugate pole is also removed\footnote{The author acknowledges the use of AI in the implementation of this step. The baseline method implementation was a forloop indexing over each pole index to check against the conjugate pole index. The Claude Opus 4.8 model was utilised to determine a succinct implementation}.
    \end{itemize}
    \item Parameter retreival: When producing the continuous time parameters, the negative frequency components corresponding to the conjugate of the poles are removed. This mask is only applied to the output parameters, and not the acutal poles of the system. Note that the phase produced from the residue associated with each pole component in this case will be cosine phase.
\end{itemize}

\subsection{Validation: simulated signals}
A validation check with several simulated signals was carried out to verify the algorithm performed as expected on signals with a predefined and known number of modes.

\begin{figure}
\centering
\includegraphics[scale=0.40]{figures/out/era-reconstruction.pdf}
\caption{Example of the ERA reconstruction of a simulated signal (left) and associated parallel filter reconstruction (right), signal parameters in [INSERT TABLE REF HERE]}
\label{fig:era-reconstruction}
\end{figure}

\begin{table}[htbp]
    \centering
    \caption{Comparison of true and ERA-estimated modal parameters}
    \label{tab:era_comparison}
    \begin{tabular}{c c c c c c c c c c c}
        \toprule
        & \multicolumn{3}{c}{Frequency (Hz)} & \multicolumn{3}{c}{$\alpha$} & \multicolumn{2}{c}{Amplitude} & \multicolumn{2}{c}{Phase (rad)} \\
        \cmidrule(lr){2-4} \cmidrule(lr){5-7} \cmidrule(lr){8-9} \cmidrule(lr){10-11}
        Mode & True & ERA & Error (\%) & True & ERA & Error (\%) & True & ERA & True & ERA \\
        \midrule
        1 & 300  & 300.11 & 0.037  & 0.026526 & 0.026516 & $-0.038$ & 1.50 & 1.50 & 0.00 & $\approx 0$ \\
        2 & 300  & 300.15 & 0.050  & 0.031831 & 0.031815 & $-0.050$ & 0.30 & 0.30 & 0.35 & 0.35 \\
        3 & 310  & 310.20 & 0.065  & 0.035938 & 0.035915 & $-0.064$ & 0.75 & 0.75 & 0.70 & 0.70 \\
        4 & 500  & 500.16 & 0.032  & 0.025465 & 0.025457 & $-0.031$ & 1.00 & 1.00 & 1.00 & 1.00 \\
        5 & 600  & 600.17 & 0.028  & 0.023873 & 0.023866 & $-0.029$ & 1.00 & 1.00 & 1.40 & 1.40 \\
        6 & 1000 & 1000.10& 0.010  & 0.015915 & 0.015913 & $-0.013$ & 1.00 & 1.00 & 1.70 & 1.70 \\
        7 & 2400 & 2390.10& $-0.413$ & 0.0073251& 0.0073249& $-0.003$ & 1.00 & 1.00 & 2.10 & 2.10 \\
        \bottomrule
    \end{tabular}
\end{table}